\reportpage{\datacontinued}{
\subheading{七、ICMP路径探测的软件证据}
命令图为终端回显已保存日志；字段图来自Wireshark打开本次原始抓包。
\ev{fig:trace-command}{trace-command}{ICMP路径探测完整命令输出（原始日志回显）}
\ev{fig:trace-probe}{trace-probe}{Traceroute首个探测请求：TTL=1、20字节零负载（帧1）}
}

\reportpage{\datacontinued}{
\subheading{路径探测报文与分组（题10--12）}
本节帧号均来自\capfile{06-traceroute-icmp}。macOS采用\code{traceroute -I}实现指导书的ICMP探测方式，完整输出见\figref{fig:trace-command}。
\qa{10. Traceroute与普通Ping报文的异同}{
本次两者请求均为Echo Request（Type 8、Code 0），目的同为\code{202.115.32.43}。普通Ping初始TTL为64，IPv4总长84，ICMP数据56字节（8字节时间戳加48字节填充）；Traceroute请求初始TTL从1递增，本次IPv4总长48，ICMP数据20字节全零，ICMP首部另占8字节。区别来自工具实现及参数，不能套用Windows固定填充字符串。帧1的字段和原始字节见\figref{fig:trace-probe}。
}
\ev{fig:trace-list}{trace-list}{路径探测完整报文列表及各跳实际响应（19帧）}
Wireshark列表中的\code{No response found}指未关联到匹配的Echo Reply，不表示没有任何ICMP差错；例如帧3仍可按引用字段手工配对帧4的Time Exceeded。
\qa{11. 初始TTL种数、每组探测数及变化}{
共使用4种初始TTL（1、2、3、4），各3个请求，共12个。\figref{fig:trace-list}结合下表按外层请求IPv4分组，不能把差错报文内嵌的TTL混入计数。所有请求Identifier（网络序）均为62861（\code{0xf58d}），Sequence依次1—12，每发一个请求递增1；同一组的TTL保持不变，下一组增1。
\par\noindent\begin{tabularx}{\linewidth}{|c|l|c|X|}\hline
初始TTL & 请求帧 & Sequence & 实际匹配响应帧\\\hline
1 &1、2、3&1—3&前两次无响应；第三次帧4\\\hline
2 &5、6、7&4—6&三次均未捕获响应\\\hline
3 &8、10、12&7—9&帧9、11、13\\\hline
4 &14、16、18&10—12&帧15、17、19\\\hline
\end{tabularx}\par
}
\qa{12. 实际响应地址、类型与原因}{
\figref{fig:trace-list}共捕获7个外层ICMP响应：\code{10.134.191.254}返回1个Type 11/Code 0；\code{202.115.39.102}返回3个Type 11/Code 0，均表示探测途中TTL耗尽；目标\code{202.115.32.43}返回3个Type 0/Code 0（Echo Reply），表示请求已到目标。没有捕获Type 3或Type 11/Code 1。嵌在差错响应中的Type 8是原请求，不另算一条外层响应。
}
}

\reportpage{\datacontinued}{
\subheading{首个差错响应与对应原请求}
帧3是第1跳的第三个探测，帧4是其返回的Time Exceeded。
\ev{fig:trace-third-probe}{trace-third-probe}{与首个Time Exceeded配对的第三个探测（帧3）}
\ev{fig:trace-error-outer}{trace-error-outer}{首个Time Exceeded的外层IPv4与ICMP字段（帧4）}
}

\reportpage{\datacontinued}{
\subheading{首个差错响应引用的原报文}
\ev{fig:trace-error}{trace-error}{第一个超时响应携带的原请求字段（帧4）}
引用的IPv4 Identification、ICMP Identifier与Sequence可与\figref{fig:trace-third-probe}的帧3核对；外层TTL则应从\figref{fig:trace-error-outer}读取。
}

\reportpage{\datacontinued}{
\subheading{差错响应的发送者、配对与TTL（题13--15）}
\qa{13. 第一个Time Exceeded的发送者与跳数}{
\figref{fig:trace-error-outer}中首次捕获的超时报文是帧4，由\code{10.134.191.254}返回，对应输出第1跳的第三个探测（\figref{fig:trace-command}）。第1跳前两个位置为星号，不能把帧4配给帧1；也不能用本机默认网关地址代替实测的响应源地址。
}
\qa{14. 所选差错响应对应哪个探测请求}{
帧4对应帧3。内嵌IPv4源/目的为\code{10.135.133.108}$\to$\code{202.115.32.43}，Identification为\code{0xf590}、总长48；引用ICMP为Type 8/Code 0、Identifier \code{0xf58d}、Sequence 3、Checksum \code{0x026f}，与帧3一致。帧4外层地址反向返回本机，时间为18:44:26.449052，帧3为18:44:26.443102（北京时间），相隔5.950毫秒。上述端点、引用字段及时间共同配对；请求见\figref{fig:trace-third-probe}，差错引用见\figref{fig:trace-error}。
}
\qa{15. 外层与内嵌IPv4的TTL分别为何}{
所选帧4外层TTL为64，属于路由器新生成并返回本机的差错IPv4报文；内层TTL为1，属于被引用的原探测IPv4报文，不能当作差错响应的TTL。引用首部可能保留丢弃前的值，不能预设必须为0。另一个路由器的帧9外层TTL为253、内层仍为1，进一步说明两层独立。外层字段见\figref{fig:trace-error-outer}，引用层见\figref{fig:trace-error}。
}
}

\reportpage{\datacontinued}{
\subheading{目标应答与到达判断（题16）}
\ev{fig:trace-final}{trace-final}{目标返回的Echo Reply（帧19）}
\qa{16. 如何判断到达目标，本次是否到达}{
ICMP方式以收到与探针Identifier和Sequence匹配、且由目标地址返回的Echo Reply判定到达。\figref{fig:trace-list}中帧14、16、18分别以TTL=4发送Sequence 10、11、12；帧15、17、19均由\code{202.115.32.43}返回Type 0/Code 0，Identifier仍为\code{0xf58d}，数据均为20字节零，分别对应三次请求，故本次已到达；帧19的展开字段见\figref{fig:trace-final}。工具完成第4跳的3次探测后停止，没有执行至最大30跳。不能套用UDP traceroute的“端口不可达”终止条件。
}
}

\reportpage{\datacontinued}{
\subheading{星号含义与证据边界（题17）}
\qa{17. 星号含义，能否据此断言后续网络不可达}{
星号表示一次探测在设定等待时间内未收到可匹配响应，本次\code{-w 1}为每次等待1秒。可能原因包括探针或回包丢失、路由器限速/过滤、回程问题等；本机证据不足以确定具体原因。第1跳前两次、第2跳三次均显示星号，但第3跳仍返回超时，第4跳仍收到目标应答（\figref{fig:trace-command}、\figref{fig:trace-list}），因此单跳星号不能证明其后网络不可达。
}
\subheading{证据边界与理论依据}
抓包只观察本机\code{en0}，未捕获响应不等于证明远端没有生成响应，也不能据本次结果确定全部路由器的限速策略。RFC 792规定Echo与Time Exceeded字段及校验范围；RFC 791规定IPv4长度、分片与偏移；RFC 1812说明差错报文引用范围及TTL处理。Smurf与畸形分片的历史风险参见RFC 4732及CVE-1999-0128。具体链接保存在附件\code{reading-notes.txt}。
}

